% Options for packages loaded elsewhere
\PassOptionsToPackage{unicode}{hyperref}
\PassOptionsToPackage{hyphens}{url}
\PassOptionsToPackage{dvipsnames,svgnames,x11names}{xcolor}
%
\documentclass[
  11pt,
]{article}
\usepackage{amsmath,amssymb}
\usepackage{iftex}
\ifPDFTeX
  \usepackage[T1]{fontenc}
  \usepackage[utf8]{inputenc}
  \usepackage{textcomp} % provide euro and other symbols
\else % if luatex or xetex
  \usepackage{unicode-math} % this also loads fontspec
  \defaultfontfeatures{Scale=MatchLowercase}
  \defaultfontfeatures[\rmfamily]{Ligatures=TeX,Scale=1}
\fi
\usepackage{lmodern}
\ifPDFTeX\else
  % xetex/luatex font selection
\fi
% Use upquote if available, for straight quotes in verbatim environments
\IfFileExists{upquote.sty}{\usepackage{upquote}}{}
\IfFileExists{microtype.sty}{% use microtype if available
  \usepackage[]{microtype}
  \UseMicrotypeSet[protrusion]{basicmath} % disable protrusion for tt fonts
}{}
\makeatletter
\@ifundefined{KOMAClassName}{% if non-KOMA class
  \IfFileExists{parskip.sty}{%
    \usepackage{parskip}
  }{% else
    \setlength{\parindent}{0pt}
    \setlength{\parskip}{6pt plus 2pt minus 1pt}}
}{% if KOMA class
  \KOMAoptions{parskip=half}}
\makeatother
\usepackage{xcolor}
\usepackage[left=1.25in,right=1.25in,top=1in,bottom=1in]{geometry}
\usepackage{longtable,booktabs,array}
\usepackage{calc} % for calculating minipage widths
% Correct order of tables after \paragraph or \subparagraph
\usepackage{etoolbox}
\makeatletter
\patchcmd\longtable{\par}{\if@noskipsec\mbox{}\fi\par}{}{}
\makeatother
% Allow footnotes in longtable head/foot
\IfFileExists{footnotehyper.sty}{\usepackage{footnotehyper}}{\usepackage{footnote}}
\makesavenoteenv{longtable}
\setlength{\emergencystretch}{3em} % prevent overfull lines
\providecommand{\tightlist}{%
  \setlength{\itemsep}{0pt}\setlength{\parskip}{0pt}}
\setcounter{secnumdepth}{-\maxdimen} % remove section numbering
\ifLuaTeX
  \usepackage{selnolig}  % disable illegal ligatures
\fi
\IfFileExists{bookmark.sty}{\usepackage{bookmark}}{\usepackage{hyperref}}
\IfFileExists{xurl.sty}{\usepackage{xurl}}{} % add URL line breaks if available
\urlstyle{same}
\hypersetup{
  pdftitle={Candidates for a first-principles derivation of (r\_e/r\_0)},
  pdfauthor={Trey Morris with Claude Opus 4.7},
  pdfsubject={Four candidate starting points for deriving r\_e from the
dual-Dirac framework's internal structure (follow-up to the
triangulation note)},
  colorlinks=true,
  linkcolor={blue},
  filecolor={Maroon},
  citecolor={Blue},
  urlcolor={blue},
  pdfcreator={LaTeX via pandoc}}

\title{Candidates for a first-principles derivation of (r\_e/r\_0)}
\author{Trey Morris with Claude Opus 4.7}
\date{2026-06-01}

\begin{document}
\maketitle

\hypertarget{candidates-for-a-first-principles-derivation-of-r_er_0-for-tepper-gill}{%
\section{\texorpdfstring{Candidates for a first-principles derivation of
\(r_e/r_0\) --- for Tepper
Gill}{Candidates for a first-principles derivation of r\_e/r\_0 --- for Tepper Gill}}\label{candidates-for-a-first-principles-derivation-of-r_er_0-for-tepper-gill}}

\textbf{Date:} 2026-05-26 (revised 2026-05-26 with two waves of
overnight progress on all three candidates). \textbf{From:} Trey Morris
(with Claude Opus 4.7). \textbf{Re:} Follow-up to the 2026-05-26
triangulation note\footnote{Triangulation note:
  \texttt{Roadmapping/Author\_Reports/2026-05\_re\_triangulation\_followup\_for\_gill.md}.};
Scope 1 of the master \emph{r\_e} tracker\footnote{Issue \#54 (Scope 1
  of the \(r_e\) tracker),
  \texttt{github.com/temoTxt/PyPhysics/issues/54}.}. The \emph{Progress
update} section reports findings from three parallel research
branches\footnote{Issue \#64 (Candidate 1, proper-time self-energy),
  \texttt{github.com/temoTxt/PyPhysics/issues/64}.}\footnote{Issue \#65
  (Candidate 2, variational route),
  \texttt{github.com/temoTxt/PyPhysics/issues/65}.}\footnote{Issue \#66
  (Candidate 3, closed-form Schwinger),
  \texttt{github.com/temoTxt/PyPhysics/issues/66}.} under their common
master\footnote{Issue \#67 (master tracker for \#64--\#66),
  \texttt{github.com/temoTxt/PyPhysics/issues/67}.}. Two of the three
branches have now reached definitive verdicts. \textbf{Candidate 2} is
halted at Outcome D: the published §III.D non-relativistic expansion is
structurally inadequate to pin \(r_e\) at any scale where the expansion
is valid. \textbf{Candidate 3} is halted at Outcome C-as-published: the
closed-form Schwinger match is algebraically forced by the back-fit
definition of the triangulated cutoff, not evidence of intentional
encoding. Candidate 1 is live, scaffolded, and on a clear next-step
plan.

\begin{center}\rule{0.5\linewidth}{0.5pt}\end{center}

\hypertarget{context}{%
\subsection{Context}\label{context}}

The DRQM I §III.D value of \(r_e\) was, originally, a cutoff parameter
fixed by a numerical search against the measured electron \(g_s\). The
empirical refinement in a recent campaign pull request\footnote{Pull
  request \#62 (empirical refinement of \(r_e/r_0\) via joint fit),
  \texttt{github.com/temoTxt/PyPhysics/pull/62}.} widened that search to
a joint fit across six observables (electron \(g_s\), hydrogen
\(2P_{3/2}-2P_{1/2}\), hydrogen \(1S\) hyperfine, helium
\({}^3P_0 - {}^3P_1\), positronium ortho-para, muonium hyperfine), and
returned \[
r_e/r_0 \;=\; 0.499\,420\,509\,912\,831\,7. \tag{1}
\] The honest reading of (1), recorded in detail in the cross-comparison
document\footnote{Cross-comparison document, §2 (``The \(r_e\)
  self-consistency across six \(g_s\)-dependent observables at the
  triangulated value''):
  \texttt{Roadmapping/Quantum\_Mechanics/Bethe\_Salpeter/10\_CrossComparison.md}.},
is the following. The six observables are not six independent
constraints on the cutoff. Each is computed as \[
\mathcal{O}_n \;=\; \bigl( g_s/-2 \bigr)^n \,\times\, \mathcal{O}_n^{\text{textbook}}, \qquad n \in \{1,2\}, \tag{2}
\] so all six share a single dependence on \(g_s\), and therefore on
\(r_e\). The joint fit is one structural fact --- the single-parameter
cutoff prescription --- applied to six manifestations, not six
independent tests. What the multi-observable fit does confirm is that
the six manifestations are mutually consistent at the framework's
precision floor. The scaling (2) holds across all six. We read this as a
self-consistency check on the framework's structural prescription, not
as an independent corroboration of the cutoff's value.

There is, at present, no first-principles derivation of \(r_e/r_0\)
anywhere in the campaign, or in the published DRQM I material we have
visibility into. The value originated as a cutoff search, and it remains
a cutoff value. The triangulation confirms that the structure (2) is
self-consistent across multiple observables under a single cutoff. It
does not promote that cutoff to a derived quantity. A genuine
first-principles derivation, in which \(r_e\) emerges from the
dual-Dirac equation's internal structure rather than being fit to data,
would be a new piece of theory work. We note that Candidate 3 below
offers a partial identification. The triangulated value already has a
clean closed-form Schwinger reading, which recasts the derivation
question into ``why does the framework specify this cutoff
prescription?'' rather than ``what numerical value should the cutoff
have?''

This note collects the three candidate starting points we have
considered for such a derivation. We give one to three paragraphs on
each, describing what the route would look like, what its plausibility
appears to be from outside the framework's internal logic, and what
would be required to pursue it. A fourth candidate we had initially
considered (retrieving a derivation step from the original DRQM I §III.D
working notebook that did not make it into the published prose) has been
ruled out by your 2026-05-25 author guidance. The original cutoff was a
numerical search against \(g_s\) alone, with no derivation-step working
draft to retrieve.

We are not proposing any of these three. We defer to your judgment on
which, if any, is natural for the framework. The thread is non-urgent.
The triangulated value already serves as the campaign's
current-best-refinement at the framework's precision floor. A
first-principles derivation is a would-be-nice-to-have, not a
load-bearing item.

\begin{center}\rule{0.5\linewidth}{0.5pt}\end{center}

\hypertarget{candidate-1-proper-time-self-energy-integral}{%
\subsection{Candidate 1 --- Proper-time self-energy
integral}\label{candidate-1-proper-time-self-energy-integral}}

The dual theory's distinguishing feature is the proper-time first-order
Dirac structure. The standard Dirac equation \[
(i\gamma^\mu \partial_\mu - m)\psi \;=\; 0 \tag{3}
\] is replaced by a proper-time formulation in which the dynamical
variable is \(\partial_\tau\) rather than \(\partial_t\). The natural
place to look for \(r_e\) as a derived quantity is the framework's
analog of the one-loop electron self-energy diagram. That is the diagram
in which a single photon propagator and a single electron propagator
close into a loop on the external electron line.

In standard QED, the one-loop self-energy \(\Sigma(p)\) carries a UV
divergence which is absorbed into mass renormalisation via a
counterterm. The cutoff procedure (Pauli--Villars, dimensional
regularisation, lattice cutoff) introduces a regulator scale, and in the
renormalised expression that scale's dependence cancels modulo the
running of \(\alpha\). In the dual theory's proper-time formulation, the
analog calculation would feature the proper-time propagator and a
proper-time photon propagator. The UV behaviour might be qualitatively
different from standard QED. It is not obvious that the integral is
divergent at all; the proper-time integral's analytic properties could
render it finite, or supply a different divergence structure.

Concretely, the route would set up a Schwinger-like proper-time integral
representation of the dual-Dirac propagator, compute the one-loop
self-energy with the proper-time photon propagator, and identify the
natural scale at which the mass counterterm equals the
physical-minus-bare mass difference. If the integral is finite, that
scale is determined unambiguously. If divergent, it is set by the
renormalisation condition. The output \(r_e/r_0\) would then be a
numerical function of \(\alpha\) and the dual-Dirac structure constants
(\(b\), \(c\), the projection-operator coefficients), evaluating to a
definite numerical value comparable against the triangulated
\(0.499\,420\,509\ldots\) of (1).

\begin{center}\rule{0.5\linewidth}{0.5pt}\end{center}

\hypertarget{candidate-2-variational-determination-via-the-renormalised-dual-dirac-equation}{%
\subsection{Candidate 2 --- Variational determination via the
renormalised dual-Dirac
equation}\label{candidate-2-variational-determination-via-the-renormalised-dual-dirac-equation}}

In the dual-Dirac equation \[
\hat{H}_{\text{dual}} \,\psi \;=\; E\,\psi, \tag{4}
\] the cutoff \(r_e\) appears as a parameter, either as a hard cutoff in
the radial integration or as a regulator in the small-\(r\) behaviour of
the wavefunction. The eigenvalue \(E\) depends on \(r_e\), so we write
\(E = E(r_e)\). A variational determination would fix \(r_e\) by
demanding that the physical electron mass eigenvalue be recovered at the
cutoff, \[
E(r_e) \;=\; m_e c^2, \tag{5}
\] including any binding-energy or self-energy contributions specified
by the framework.

This is analogous to how the standard-QED running coupling
\(\alpha(\mu)\) is fixed by the renormalisation condition
\(\alpha(\mu_0) = \alpha_{\text{phys}}\) at a chosen reference scale. A
framework parameter is fixed by a physical observable evaluated at the
framework's natural scale. The advantage of this route over Candidate 1
is that it side-steps the UV-divergence question entirely. The
variational principle operates within the framework's renormalised
structure, not on the divergent bare-vertex calculation. The
disadvantage is that the mass condition (5) alone may not uniquely
determine \(r_e\). Additional consistency conditions (gauge invariance
of the photon propagator at the cutoff, current conservation at the
radial boundary, perhaps the magnetic-moment relation closing the
system) may be needed for closure.

If the right set of consistency conditions can be assembled, the
variational route is the most direct path from the dual-Dirac equation's
internal structure to a numerical \(r_e/r_0\). The calculation would be
a coupled eigenvalue and renormalisation-condition problem, solvable
symbolically once the conditions are specified. The output is again
comparable against the triangulated value of (1). The route's
credibility, however, depends on whether the additional consistency
conditions are framework-internally motivated or look ad hoc.

\begin{center}\rule{0.5\linewidth}{0.5pt}\end{center}

\hypertarget{candidate-3-the-triangulated-value-already-has-a-clean-closed-form-schwinger-reading}{%
\subsection{Candidate 3 --- The triangulated value already has a clean
closed-form Schwinger
reading}\label{candidate-3-the-triangulated-value-already-has-a-clean-closed-form-schwinger-reading}}

We observe that the triangulated value (1) is, to within
\(\sim 10^{-6}\), the closed-form value \[
r_e/r_0 \;=\; \frac{2 - \alpha/(2\pi)}{4 + \alpha/\pi} \;=\; 0.499\,419\,632\,156\ldots, \tag{6}
\] obtained by inverting \[
g_r \;=\; 2\Bigl[\,1 \,-\, \frac{4\,r_0}{2r + r_0}\,\Bigr] \tag{7}
\] against the Schwinger one-loop QED anomalous moment
\(g_e^{(1\text{-loop})} = -2 - \alpha/\pi = -2.002\,322\,819\,465\,7\ldots\)
The residual discrepancy between the closed-form \(0.499\,419\,632\) of
(6) and the triangulated \(0.499\,420\,510\) of (1) is exactly the gap
between the Schwinger one-loop and the all-orders measured
\(g_e = -2.002\,319\,304\,362\,5\ldots\) The latter includes the
Karplus--Kroll two-loop and higher corrections at \(\sim 10^{-6}\)
relative. In other words, at one-loop QED precision, the cutoff
prescription \[
r_e \;=\; \frac{2 - \alpha/(2\pi)}{4 + \alpha/\pi}\,r_0 \tag{8}
\] is what makes (7) reproduce the Schwinger result exactly.

This is, from outside the framework's internal logic, the cleanest and
most natural reading of the cutoff identification. The DRQM I §III.D
formula is, at the precision the framework's apparatus delivers, an
algebraic re-encoding of the Schwinger one-loop anomalous moment through
a particular cutoff substitution. The reason the triangulation across
six observables returns essentially the same value as the uni-observable
back-fit against \(g_s\) alone is that every other observable's
prediction follows the scaling (2). Substituting the measured \(g_s\) is
the same operation whether done once or six times. The measured \(g_s\)
is, to one-loop QED precision, exactly what the closed-form cutoff (8)
is engineered to reproduce.

On this reading, the first-principles derivation question reduces to:
why does the framework specify this particular cutoff prescription? Two
sub-questions follow. First, is there a derivation in which the
dual-Dirac equation's renormalisation structure produces (8) as the
natural mass-renormalisation scale (Candidates 1 or 2 applied at
one-loop precision)? Second, is the cutoff a structural constant of the
dual representation itself, such that the combination
\((2 - \alpha/(2\pi))/(4 + \alpha/\pi)\) comes out of the framework's
algebra without reference to QED at all? The first is a one-loop
calculation that should be tractable if the framework's renormalisation
prescription is fully specified. The second requires identifying the
framework-internal origin of an \(\alpha\)-dependent ratio, which is
harder. If the first reading is the natural one, the campaign can
attempt the calculation directly. If the second, we would need your
guidance on which structural decomposition produces \(\alpha\)-dependent
ratios in the framework.

We did not see this identification in the published DRQM I §III.D text.
If it is a known property of the framework, please tell us; that would
simplify the disposition considerably. If it is a coincidence at the
campaign's precision (the triangulated value happens to be close to the
Schwinger closed-form but is not engineered to be), the
Karplus--Kroll-level residual is the test. With more
precision-spectroscopy observables in the joint fit, the residual should
not systematically point at the two-loop QED corrections unless the
cutoff genuinely is encoding QED.

\begin{center}\rule{0.5\linewidth}{0.5pt}\end{center}

\hypertarget{progress-update-2026-05-26-pm-three-parallel-research-branches}{%
\subsection{Progress update (2026-05-26 PM) --- three parallel research
branches}\label{progress-update-2026-05-26-pm-three-parallel-research-branches}}

Since the morning version of this note went out, three parallel research
branches\footnote{Issue \#64 (Candidate 1, proper-time self-energy),
  \texttt{github.com/temoTxt/PyPhysics/issues/64}.}\footnote{Issue \#65
  (Candidate 2, variational route),
  \texttt{github.com/temoTxt/PyPhysics/issues/65}.}\footnote{Issue \#66
  (Candidate 3, closed-form Schwinger),
  \texttt{github.com/temoTxt/PyPhysics/issues/66}.} under their common
master\footnote{Issue \#67 (master tracker for \#64--\#66),
  \texttt{github.com/temoTxt/PyPhysics/issues/67}.} have begun
attempting the three candidates. Per Crocco compliance, the work below
is substantive AI, with per-paragraph human-review markers in the
per-branch state logs\footnote{Per-branch state log on each Candidate's
  branch: \texttt{.dev/research/STATE.md}.}. The headline findings are
summarised here for your morning read.

\hypertarget{candidate-1-proper-time-self-energy3-live-conceptual-re-framing-of-r_e-baseline-scaffolded}{%
\subsubsection[Candidate 1 (proper-time self-energy) --- live;
conceptual re-framing of \(r_e\), baseline
scaffolded]{\texorpdfstring{Candidate 1 (proper-time
self-energy)\footnote{Issue \#64 (Candidate 1, proper-time self-energy),
  \texttt{github.com/temoTxt/PyPhysics/issues/64}.} --- live; conceptual
re-framing of \(r_e\), baseline
scaffolded}{Candidate 1 (proper-time self-energy) --- live; conceptual re-framing of r\_e, baseline scaffolded}}\label{candidate-1-proper-time-self-energy3-live-conceptual-re-framing-of-r_e-baseline-scaffolded}}

Four iterations: read DRQM I §II--III, \emph{The Classical Electron
Problem}, and the Bethe--Salpeter Lamb-shift document; scaffolded the
dedicated notebook\footnote{Mathematica notebook for the Candidate 1
  derivation:
  \texttt{Roadmapping/Mathematica\_Notebooks/Quantum\_Mechanics/r\_e\_derivation\_self\_energy.wl}.}
with a baseline cell. The baseline standard-QED on-shell shift, \[
\frac{\delta m}{m} \;=\; \frac{3\alpha}{4\pi}\,\Bigl[\,\log\!\frac{\Lambda^2}{m^2} \,+\, \tfrac{1}{2}\,\Bigr], \tag{9}
\] is the Bjorken--Drell Eq. (10.59) result, verified symbolically.

The headline finding is conceptual rather than computational. The
existing Lamb-shift calculation uses the textbook non-relativistic Bethe
(1947) UV cutoff \(K \sim m c^2\), equivalently
\(\lambda_C = \hbar/(mc)\), while the triangulated
\(r_e \sim r_0/2 = (\alpha/2)\,\lambda_C\). Thus \(r_e\) is
parametrically smaller than the Bethe cutoff by \[
r_e / \lambda_C \;\sim\; \alpha/2 \;\approx\; 3.65 \times 10^{-3}. \tag{10}
\] We conclude that \(r_e\) is not a UV-loop cutoff at all. It sits at
the Coulomb-binding scale, not the Compton scale. Indeed, at
\(r = r_0/2\), \[
\lvert V_0 \rvert \;=\; 2\,m c^2, \tag{11}
\] which is exactly the pair-production threshold. The natural
re-reading is that \(r_e\) marks the radius inside which the bound-state
wavefunction picks up virtual \(e^+ e^-\) contributions, and inside
which the §III.D small-component elimination \[
\psi_2 \;=\; \frac{c\,(\boldsymbol{\sigma}\!\cdot\!\boldsymbol{\pi})\,\psi_1}{\lambda - V_0 + m c^2} \tag{12}
\] stops being a valid approximation.

Technically, the (II.3) ``potential-in-the-mass'' form is the cleanest
kernel. There are no
\(V \boldsymbol{\alpha}\!\cdot\!\boldsymbol{\pi}/(mc)\) or
\(\boldsymbol{\alpha}\!\cdot\!\boldsymbol{\nabla} V/(mc)\) pieces (those
arise from operator non-commutativity in the (II.1) Dirac form), so the
kernel reduces to the Pauli kinetic plus \(V^2/(2 m c^2)\).

A heuristic sanity check at iteration 4 sharpens the discriminator. The
naive identification \(\Lambda = \hbar/(r_e c)\) at \(r_e/r_0 = 0.5\)
gives \[
\log(\Lambda^2/m^2) \;=\; 11.23, \qquad \delta m/m \;=\; 2.04 \times 10^{-2}, \tag{13}
\] versus the natural one-loop coupling
\(\alpha/(4\pi) = 5.81 \times 10^{-4}\) --- a \(35\times\) overestimate.
Two distinguishable resolutions are queued for the next cells. First,
the framework may supply \(\Lambda \sim \hbar/(b\,r_e)\) with \(b > c\)
for bound states, suppressing the log. Second, the Bethe-log replacement
\(\log(K/\langle\Delta E\rangle)\) at fixed photon-loop measure
(sum-over-states log) is parametrically smaller. The two are
distinguishable. The first shifts \(r_e/r_0\) via \(\log(b/c)\) at fixed
\(\langle p^2\rangle\); the second shifts \(r_e/r_0\) via the Bethe-log
replacement at fixed photon-loop measure. \textbf{Status: live, on
plan.} The outcome-matrix branch remains open. The next two cells are
expected to discriminate the two resolutions and produce a numerical
\(r_e/r_0\).

\hypertarget{candidate-2-variational-route4-halted-at-outcome-d-the-published-nr-expansion-is-structurally-inadequate-to-pin-r_e}{%
\subsubsection[Candidate 2 (variational route) --- halted at Outcome D;
the published NR-expansion is structurally inadequate to pin
\(r_e\)]{\texorpdfstring{Candidate 2 (variational route)\footnote{Issue
  \#65 (Candidate 2, variational route),
  \texttt{github.com/temoTxt/PyPhysics/issues/65}.} --- halted at
Outcome D; the published NR-expansion is structurally inadequate to pin
\(r_e\)}{Candidate 2 (variational route) --- halted at Outcome D; the published NR-expansion is structurally inadequate to pin r\_e}}\label{candidate-2-variational-route4-halted-at-outcome-d-the-published-nr-expansion-is-structurally-inadequate-to-pin-r_e}}

Four iterations enumerated seven closure conditions. Only the
mass-renormalisation condition (closure \#7) was both framework-internal
and potentially of sufficient precision. We pursued it under the
textbook working hypothesis
\(\Delta E_{\rm bind}^{\rm framework} = \langle V_0\rangle\),
\(\Delta E_{\rm SE}^{\rm framework} = 0\), with trial
\(\psi_1 = N e^{-r/aa}\) on the cutoff-restricted domain
\([r_e, \infty)\). The closure equation, in the dimensionless variables
\(\hat{a} = aa/r_0\) and \(\hat{r}_e = r_e/r_0\), is \[
E_{\rm dim}(\hat{a}, \hat{r}_e) \;=\; \frac{1}{2\alpha^2 \hat{a}^2} \;-\; \frac{\hat{a} + 2\hat{r}_e - 1}{\hat{a}^2 + 2\hat{a}\hat{r}_e + 2\hat{r}_e^2} \;=\; 0. \tag{14}
\]

The diagnostic table at \(\alpha = 1/137.035999\) exposes the structural
inadequacy.

\begin{longtable}[]{@{}
  >{\raggedright\arraybackslash}p{(\columnwidth - 8\tabcolsep) * \real{0.2000}}
  >{\raggedright\arraybackslash}p{(\columnwidth - 8\tabcolsep) * \real{0.2000}}
  >{\raggedright\arraybackslash}p{(\columnwidth - 8\tabcolsep) * \real{0.2000}}
  >{\raggedright\arraybackslash}p{(\columnwidth - 8\tabcolsep) * \real{0.2000}}
  >{\raggedright\arraybackslash}p{(\columnwidth - 8\tabcolsep) * \real{0.2000}}@{}}
\toprule\noalign{}
\begin{minipage}[b]{\linewidth}\raggedright
Regime
\end{minipage} & \begin{minipage}[b]{\linewidth}\raggedright
\(\hat{a}\)
\end{minipage} & \begin{minipage}[b]{\linewidth}\raggedright
\(\hat{r}_e\)
\end{minipage} & \begin{minipage}[b]{\linewidth}\raggedright
\(E_{\rm dim}\)
\end{minipage} & \begin{minipage}[b]{\linewidth}\raggedright
Verdict
\end{minipage} \\
\midrule\noalign{}
\endhead
\bottomrule\noalign{}
\endlastfoot
Electron-radius scale & \(1\) & \(0.5\) & \(+9389.0\) & Kinetic
dominates by four orders; \textbf{NR expansion invalid}. \\
Bohr scale & \(1/\alpha^2 \approx 18\,778\) & \(0.5\) &
\(-2.662\times 10^{-5}\) & Matches \(-\alpha^2/2\); \textbf{cutoff
invisible}. \\
Bohr scale, no cutoff & \(1/\alpha^2\) & \(0\) &
\(-2.662\times 10^{-5}\) & Identical to \(\hat{r}_e = 0.5\) to ten
significant figures. \\
\end{longtable}

At the cutoff scale, both expansion parameters of (III.4), namely
\(V_0/(m c^2)\) and \(\hbar/(m c\,r)\), are \(O(1)\). The expansion is
invalid there. At the Bohr scale, where the expansion is valid, the
trial mass-density \(r^2 e^{-2r/aa}\) peaks at
\(r \sim aa \sim 18\,000\,r_0\), suppressing the cutoff coupling by
\(\alpha^2\). It follows that no scale exists at which both (a) the NR
expansion is valid and (b) the cutoff couples meaningfully to the
closure. The published expanded \(K_D\) cannot pin \(r_e\).
\textbf{Status: halted, Outcome D, blocked.} Two specific clarifications
would unblock. First, a framework-internal
\(\Delta E_{\rm SE}^{\rm framework}(r_e)\) at the cutoff scale. Second,
a redirect to the un-expanded full \(H_D\) under a radial-cutoff
regulator, which would be a five- to ten-iteration arc, distinct from
any previously enumerated route, and not initiated without your
endorsement. See Question 2 below.

\hypertarget{candidate-3-closed-form-schwinger5-halted-at-outcome-c-as-published-the-kk-lr-kf-residual-match-is-algebraically-forced}{%
\subsubsection[Candidate 3 (closed-form Schwinger) --- halted at Outcome
C-as-published; the KK + LR + KF residual match is algebraically
forced]{\texorpdfstring{Candidate 3 (closed-form Schwinger)\footnote{Issue
  \#66 (Candidate 3, closed-form Schwinger),
  \texttt{github.com/temoTxt/PyPhysics/issues/66}.} --- halted at
Outcome C-as-published; the KK + LR + KF residual match is algebraically
forced}{Candidate 3 (closed-form Schwinger) --- halted at Outcome C-as-published; the KK + LR + KF residual match is algebraically forced}}\label{candidate-3-closed-form-schwinger5-halted-at-outcome-c-as-published-the-kk-lr-kf-residual-match-is-algebraically-forced}}

Four iterations executed the empirical-test path. At twenty-digit
precision with \(\alpha = 7.297\,352\,569\,3 \times 10^{-3}\), we record
the comparison.

\begin{longtable}[]{@{}
  >{\raggedright\arraybackslash}p{(\columnwidth - 2\tabcolsep) * \real{0.5000}}
  >{\raggedright\arraybackslash}p{(\columnwidth - 2\tabcolsep) * \real{0.5000}}@{}}
\toprule\noalign{}
\begin{minipage}[b]{\linewidth}\raggedright
Quantity
\end{minipage} & \begin{minipage}[b]{\linewidth}\raggedright
Value
\end{minipage} \\
\midrule\noalign{}
\endhead
\bottomrule\noalign{}
\endlastfoot
\(r_e/r_0\) closed-form \((2-\alpha/(2\pi))/(4+\alpha/\pi)\) &
\(0.499\,419\,632\,156\,99\) \\
\(r_e/r_0\) triangulated (PR-refined) & \(0.499\,420\,509\,912\,83\) \\
\(\Delta g_e^{\rm obs} = g_e^{\rm meas} - g_e^{\rm Schwinger}\) &
\(+3.5151 \times 10^{-6}\) \\
Karplus--Kroll two-loop \(+2 C_2 (\alpha/\pi)^2\),
\(C_2 = 0.328\,478\,965\,579\) & \(+3.5446 \times 10^{-6}\) \\
Laporta--Remiddi three-loop \(-2 C_3 (\alpha/\pi)^3\),
\(C_3 = 1.181\,241\,456\,587\) & \(-2.961 \times 10^{-8}\) \\
Kinoshita--Fukuda four-loop \(+2 C_4 (\alpha/\pi)^4\),
\(C_4 \approx 1.9106\) & \(+1.1 \times 10^{-10}\) \\
Sum (KK + LR + KF) & \(+3.5151 \times 10^{-6}\) \\
Residual: \(\Delta g_e^{\rm obs} - \Delta g_e^{\rm pred,\,all}\) &
\(-9.7 \times 10^{-12}\) \\
\end{longtable}

Iteration 3 read DRQM I §III.D Eqs. (III.18)--(III.23) line by line and
confirmed that the as-published §III.D derivation does not produce
\(r_e/r_0\) as a closed form in \(\alpha\). The numerical value at
(III.22), \(r_e = 0.499\,857\,150\,068\,631\,r_0\), is presented as an
empirical fact (the uni-observable numerical search you confirmed on
2026-05-25), not derived. Consequently, the \(10^{-11}\) agreement above
is algebraically forced by the back-fit definition of the triangulated
\(r_e\). Any cutoff satisfying \(g_r(r) = g_e^{\rm meas}\) will, when
subtracted from the closed-form one-loop value, yield via \(dg_r/dr\)
propagation a \(\Delta g_e\) identical to the
all-orders-QED-beyond-one-loop content of measured \(g_e\). That is, it
must yield KK + LR + KF. The KK-consistency observation is necessary but
not sufficient for intentional encoding. The as-published apparatus does
not supply the sufficient piece.

\textbf{Status: halted, Outcome C-as-published. Finding 2 stays Marginal
(characterised but not promoted).} The canonical record is now committed
as a new subsection of the DRQM I verification document\footnote{DRQM I
  verification document, §III.D subsection (``Schwinger identification
  --- empirical residual test''):
  \texttt{Roadmapping/Equation\_Verification/Dual\_Relativistic\_Quantum\_Mechanics\_I.md}.}
under §III.D (``Schwinger identification --- empirical residual test'').
The path to Outcome B (intentional Schwinger encoding implies Finding 2
promoted to \textbf{Pass}) now passes solely through a first-principles
rederivation that produces the closed-form as a derived identity, that
is, through Candidate 1 if it lands on the closed form, or through a
sub-route of Candidate 2 if author input redirects to the un-expanded
full \(H_D\).

\hypertarget{consolidated-questions-for-you}{%
\subsubsection{Consolidated questions for
you}\label{consolidated-questions-for-you}}

After this overnight run, two questions remain load-bearing (the third
was resolved by the iter-3 §III.D line-by-line read). Even one answer
accelerates the remaining live branch.

\begin{enumerate}
\def\labelenumi{\arabic{enumi}.}
\item
  \textbf{(Candidate 1, live branch.)} Does the framework specify a
  proper-time photon propagator form? Two natural candidates surface.
  First, a Schwinger proper-time form with \(b\) replacing \(c\) in the
  dispersion, \[
  k^2 \;=\; (\omega/b)^2 \,-\, \mathbf{k}^2. \tag{15}
  \] Second, the standard Feynman propagator with the \(b/c\) conversion
  absorbed into the source coupling rather than the propagator. The two
  give different numerical \(r_e/r_0\) at the same order in \(\alpha\).
  The iter-4 sanity check (13) shows the difference is at the
  \(\log(b/c)\) level, and hence discriminable.
\item
  \textbf{(Candidate 2, halted, unblocking.)} The published expanded
  \(K_D\) of (III.4) is structurally inadequate to pin \(r_e\), as the
  diagnostic table above shows. Two unblocking moves present themselves.
  \textbf{(2a)} Does the framework specify a
  \(\Delta E_{\rm SE}^{\rm framework}(r_e)\) at the cutoff scale that we
  should have been carrying in
  \(\langle K_D \rangle = m_e c^2 + \Delta E_{\rm bind} + \Delta E_{\rm SE}\)?
  \textbf{(2b)} Is the variational determination intended to operate on
  the un-expanded full \(H_D\) under a radial-cutoff regulator
  (radial-Dirac eigenvalue problem on \(r \in [r_e, \infty)\)) rather
  than the expanded \(K_D\) of (III.4)? The un-expanded route is a five-
  to ten-iteration arc we have not begun without your endorsement.
\end{enumerate}

The original Question 3 from earlier today (does §III.D derive
\(r_e/r_0\) as a closed form in \(\alpha\)?) was answered by the iter-3
line-by-line read. As published, §III.D does not derive a closed form;
the numerical value is an empirical input from your uni-observable
numerical search. The remaining open piece on the closed-form-encoding
question is whether an in-progress or planned rederivation (per
Candidate 1, or a redirected Candidate 2) would produce the closed form.
That is now the substantive thread for the master tracker\footnote{Issue
  \#67 (master tracker for \#64--\#66),
  \texttt{github.com/temoTxt/PyPhysics/issues/67}.}.

\begin{center}\rule{0.5\linewidth}{0.5pt}\end{center}

\hypertarget{closing-how-this-thread-sits}{%
\subsection{Closing --- how this thread
sits}\label{closing-how-this-thread-sits}}

We expect the most likely outcome, given your 2026-05-25 guidance and
the campaign's posture, is that none of these three candidates rises
above the triangulated cutoff in usefulness for the campaign's current
scope. The triangulated value is empirically well-constrained at the
framework's precision floor (the joint-fit consistency across six
manifestations of the scaling (2)), agrees with measurement at that
floor, and is sufficient for the campaign's downstream work without a
first-principles derivation behind it. The Scope 1 thread is therefore
non-urgent. It remains open as a would-be-nice-to-have for future
framework development, not as a load-bearing item. Candidate 3's
Schwinger closed-form identification, if intentional in the framework's
construction, may already be the answer.

If one of the three candidates does look natural to you, we would be
glad to pursue it via Mathematica symbolic calculation with your
guidance on the framework's internal logic. If none of them looks
natural, the triangulated value (1) stands as the campaign's \(r_e\)
disposition and Scope 1 can be closed without a derivation. Either way,
our position is that the triangulation has confirmed the structure (2)
is self-consistent under a single cutoff at the framework's precision
floor; that is the campaign's current honest-scope position on the
\(r_e\) question.

The Progress update above adds three substantive pieces of physical
content on top of the morning version. First, the re-framing of \(r_e\)
from ``UV cutoff'' to ``Coulomb-binding scale at the pair-production
threshold'' (Candidate 1) --- a statement about what \(r_e\) is in the
framework's small-component-elimination apparatus, independent of which
derivation route eventually settles its first-principles status, and
worth your read even if Scope 1 closes without a derivation. Second, the
structural inadequacy of the published expanded \(K_D\) for variational
\(r_e\) determination (Candidate 2, Outcome D). At no scale do (a) the
validity of the (III.4) NR expansion and (b) meaningful coupling of the
cutoff to the closure equation coexist. Third, the algebraically-forced
character of the Schwinger closed-form agreement (Candidate 3, Outcome
C-as-published). The \(10^{-11}\) KK + LR + KF residual match is
necessary but not sufficient for intentional encoding, and the
as-published §III.D does not provide the sufficient piece. The
disposition of Scope 1 now turns on the live Candidate 1 branch (and on
whether Candidate 2 redirects to the un-expanded full-\(H_D\) arc with
your endorsement).

--- Trey

\begin{center}\rule{0.5\linewidth}{0.5pt}\end{center}

\hypertarget{references}{%
\subsection{References}\label{references}}

\end{document}
